\chapter{各章习题参考答案与提示}

本附录给出各章 `problemset` 习题的参考答案。需要“写代码”的题目只给出核心实现思路，读者可在完整示例工程中找到可直接运行的版本；标为“思考题”的题目没有唯一答案，这里只提供分析角度。

% ---------------------------------------------------------------------
\section{第~\ref{ch:embedded-c}~章 参考答案}

\begin{enumerate}
\item 十进制 173 的二进制为 \texttt{10101101}，八进制为 \texttt{255}，十六进制为 \texttt{AD}。
      8 位补码中 \texttt{-5} 为 \texttt{11111011}。
\item 置位：\texttt{REG |= (1u << 3);}；清零：\texttt{REG \&= \textasciitilde(1u << 3);}
\item 位域适合“多个有名字的小字段”反复读写，可读性好，但位分配顺序由编译器决定；联合体适合“同一块内存不同时刻只保存一种类型”的场景，通常需要额外的类型字段区分当前有效成员。
\item 大端从低地址到高地址为 \texttt{12 34 56 78}，小端为 \texttt{78 56 34 12}。可通过联合体读首字节或用 \texttt{uint8\_t *} 强转指针判断。
\item 函数指针保存的是代码入口地址，且签名必须匹配。示例：
\begin{lstlisting}[style=c]
typedef struct { void (*Init)(void); int (*Read)(void); } Device_T;
\end{lstlisting}
\item \texttt{volatile} 只阻止“访问被优化掉”，不能保证“读-改-写”原子性，也不能约束执行顺序；并发修改共享状态时仍需要关中断或临界区。
\end{enumerate}

% ---------------------------------------------------------------------
\section{第~\ref{ch:scheduler-and-state-machine}~章 参考答案}

\begin{enumerate}
\item 定时器中断负责记录时间与“置唤醒标志”，主循环负责扫描任务并执行回调；中断只做短小工作，避免阻塞其他中断。
\item TCB 至少保存状态、回调指针、下次运行时间、周期和运行标志；增加优先级时可加入 \texttt{priority} 字段，并让就绪队列按优先级排序。
\item 按键状态可拆为：空闲、按下消抖确认、按住、松开确认。消抖阶段只作状态迁移，不立即触发业务，比单纯延时更省 CPU 且能区分短按/长按。
\item 中断与主循环同时修改任务标志会形成数据竞争。\texttt{volatile} 只保证读到最新内存值，写-改-写仍要临时关中断或使用临界区。
\item 就绪队列避免每次主循环都扫描全部任务；本实现用 \texttt{isInMsgQuene} 标志防止同一任务重复入队。
\end{enumerate}

% ---------------------------------------------------------------------
\section{第~\ref{ch:simple-task}~章 参考答案}

\begin{enumerate}
\item \texttt{Task\_Tick} 在中断里推进时间并把到期的任务放入 FIFO；\texttt{Task\_Loop} 在主循环中取出任务并调用任务运行函数 \texttt{Task\_RunTask} 执行。
\item 牺牲一个槽位后，\texttt{头 == 尾} 表示空，\texttt{(尾+1)\%容量 == 头} 表示满，判断无需额外字段；若额外记录长度，生产者与消费者都要更新长度，临界区逻辑更复杂。
\item \texttt{Task\_DelayUntil} 适合固定周期任务，误差不会累积；\texttt{Task\_DelayInterrupt} 适合需要“时间到后立即被唤醒”的定时语义，注意回调不能阻塞。
\item \texttt{volatile} 解决“编译器看不到值会变”的问题；但它不保证原子性。涉及入队、出队等读-改-写时必须关闭定时器中断。
\item 静态任务池在编译期分配，无碎片、无内存耗尽风险，失败模式确定；代价是任务数量上限固定。
\item FIFO 按入队顺序执行同一批到期任务，先到先服务，因此不会出现某个任务长期得不到执行的饿死问题。
\end{enumerate}

% ---------------------------------------------------------------------
\section{第~\ref{ch:serial-fifo}~章 参考答案}

\begin{enumerate}
\item 固定数组被写满后，后续字节会覆盖未读数据或直接被丢弃，例如高波特率下主循环稍慢就丢帧。
\item 三种情况分别为：\texttt{StartIndex != EndIndex}、写入后 \texttt{EndIndex} 追上 \texttt{StartIndex} 的前一个槽位、以及写入区间跨过数组末尾后 \texttt{EndIndex} 回到头部附近。
\item 回绕后 \texttt{EndIndex - StartIndex} 可能为负；先加 \texttt{Size} 再取模可得到正确长度。
\item 入队/出队属于读-改-写，需要临时关闭对应串口中断或用临界区，否则可能同时修改头尾指针。
\item 校验码用于发现传输错误；校验失败应丢弃整帧并从“帧头”重新同步，避免把坏数据当作命令执行。
\item 逐字节中断实时性好但中断频繁；空闲中断适合不定长接收；DMA 搬运占用 CPU 最少，但代码与内存复杂度更高。
\end{enumerate}

% ---------------------------------------------------------------------
\section{第~\ref{ch:oop-in-c}~章 参考答案}

\begin{enumerate}
\item 封装的核心是“只暴露稳定接口、隐藏内部实现”，在 C 中通常用不透明结构体指针加模块函数实现。
\item 子结构体把父结构体放在第一个成员，即可通过强转共享父类布局；访问子类字段需还原为子结构体指针。
\item 函数指针表相当于手工虚函数表；新增设备只需实现同一组接口并注册，调用方不需要 \texttt{if/else} 区分类型。
\item 联合体成员共享同一块内存，未写入的成员没有确定值；必须用类型枚举先判断再读取。
\item 面向接口编程把“策略”与“调用方”解耦，例如菜单回调、存储设备读写、显示接口都可用同一套表驱动代码复用。
\end{enumerate}

% ---------------------------------------------------------------------
\section{第~\ref{ch:memory-alloc}~章 参考答案}

\begin{enumerate}
\item 典型布局从低地址到高地址依次是：.data/.bss、堆向上增长、栈向下增长；栈与堆在中间碰撞会互相覆盖导致系统崩溃。
\item \texttt{malloc/free} 时间不确定、易产生碎片、失败时缺少明确恢复路径；固定块池把所有块大小统一，释放后可直接复用，不会产生外部碎片。
\item 可给每块分配一位：\texttt{0} 表示空闲、\texttt{1} 表示占用；扫描字/字节中的位即可快速找到空块。
\item Cortex-M 的非对齐 \texttt{LDR/STR} 可能触发 UsageFault 并升级为 HardFault；因此结构体大小与指针地址都要按成员最大对齐要求处理。
\item 4 字节拷贝利用机器字长，搬运同样字节数所需循环次数更少；源与目标重叠时应改用类似 \texttt{memmove} 的逐字节安全拷贝。
\item 用多个不同块大小的池，申请时按大小选择最接近的池，既能覆盖多种尺寸又不会浪费整块内存。
\end{enumerate}

% ---------------------------------------------------------------------
\section{第~\ref{ch:param-store}~章 参考答案}

\begin{enumerate}
\item 魔数用于识别“这一块数据是否属于本参数结构体”，能发现写错版本或误读其他存储区；校验和只能发现位错误。
\item 常见策略：读取后按版本号补默认字段再写回，或使用 \texttt{schema\_v1/schema\_v2} 迁移工具生成“旧结构体到新结构体”的转换代码。
\item 分组保存可只写发生变化的参数组，减少 Flash 擦写次数，也能隔离“某组损坏”对整块参数的影响；代价是需要在每组带独立校验与版本信息。
\item 文件系统便于电脑导出/导入；二进制体积小、解析快；\texttt{json/yaml} 可读性好但解析开销大，适合 PC 工具或配置量小的场合。
\item Python 脚本、第三方库和文本处理能力强，适合写“读旧文件→转换→生成 C 迁移函数”的一次性工具，且不必占用单片机 Flash。
\end{enumerate}

% ---------------------------------------------------------------------
\section{第~\ref{ch:menu-system}~章 参考答案}

\begin{enumerate}
\item 菜单项中有一个“项类型”枚举（跳转/函数/数值/单选/多选），运行时先判断类型，再读取联合体对应成员。
\item 静态创建把内存峰值限制在编译期，避免运行期碎片；若分配越界，应触发断言或返回错误码，绝不能静默覆盖相邻内存。
\item 页面栈适合“进入子页面再返回”的层级结构，需要记录栈深；显式跳转只需保存当前页面入口，返回目标由调用方显式给出，更适合平级页面。
\item 不填充背景直接绘制文字会残留上一帧的笔画；先清区域再绘制才能保证字符不重叠。
\item 把 LCD 底层 API 隔离到移植层后，菜单核心在 PC 与单片机上使用同一份代码，可先在 PC 验证逻辑再搬硬件。
\end{enumerate}

% ---------------------------------------------------------------------
\section{第~\ref{ch:command-parser}~章 参考答案}

\begin{enumerate}
\item 切词时可用引号或转义表示字符串参数；例如把 \texttt{"hello world"} 作为一个整体 token。
\item 哈希表把字符串比较次数从“逐个命令”降为一次哈希定位；冲突采用线性探测，把后续元素放到下一个空槽。
\item 命令注册表把名字、函数与帮助集中存放，新增命令只需注册一次，主循环和帮助命令都不需要改。
\item \texttt{Param\_T} 内部用类型字段加联合体；\texttt{INT\_PARAM(param, i)} 之类宏先校验类型再按 int 读取。
\item 纯 C 没有反射，无法在运行期按名字自动匹配参数类型；宏可按“类型名展开代码”，C++ 模板可自动生成适配器，本质都是把类型信息放到编译期。
\end{enumerate}

% ---------------------------------------------------------------------
\section{第~\ref{ch:script-engine}~章 参考答案}

\begin{enumerate}
\item 终结符是不可再展开的最小 token；非终结符可被产生式展开；产生式描述展开规则；开始符号是整棵推导树的根。
\item 递归下降用“函数嵌套深度”表达优先级：乘法因子放在 \texttt{ParseFactor}，加减法在上层 \texttt{ParseExpr}，于是 \texttt{*} 先于 \texttt{+} 结合。
\item 直接左递归会让函数无限调用自身；生成器把左递归拆成“先解析基产生式，再循环匹配续产生式”，并在每轮把上一轮结果当作新 \texttt{\$1}。
\item First 集合决定“遇到某个 token 该选哪条产生式”，Follow 集合决定可空产生式何时匹配空；发生冲突时可左因子提取、改写文法或增加向前看。
\item $LR(1)$ 能力最强但表大；$LALR(1)$ 是体积与能力的折中；$LL(1)$ 实现最简单。单片机小资源下通常选择递归下降 $LL(1)$ 或手工解析器。
\item .gram 用逐条 \texttt{lhs := ...} 免去 JSON 嵌套，\texttt{\%var} 可直接声明数组/结构体/字典，动作里写 \texttt{\$1 + \$2} 和 \texttt{int32(0)} 也更接近自然语言。
\end{enumerate}

% ---------------------------------------------------------------------
\section{第~\ref{ch:coroutine-scheduler}~章 参考答案}

\begin{enumerate}
\item \texttt{setjmp} 保存程序计数器、栈指针和 callee-saved 寄存器等现场；\texttt{longjmp} 恢复后，\texttt{setjmp} 之后被修改的“非易失”局部变量值不再保证，应避免依赖。
\item 独立栈保证每个任务挂起时的局部变量都留在自己的栈中；共享栈切换前必须把“本次实际使用过的栈区域”备份，否则会被下一个任务覆盖。
\item 典型迁移：READY 由调度器选中进入 RUNNING；任务调用 \texttt{Task\_Wait} 时保存现场并进入 SUSPEND；时间到或条件满足后被重新置为 READY。
\item 协程把执行到哪一步“藏”在 \texttt{jmp\_buf} 和栈里，代码更接近顺序思维；调试时则难以直接观察“当前在第几行”。
\item 可在栈底填充固定水印值，任务切换或结束时检查水印是否被覆盖，从而得到该任务的最大栈用量。
\item MSP 是主堆栈指针，PSP 是进程堆栈指针；切换任务栈时一般修改 PSP，并在异常返回时由硬件选择对应栈指针。
\end{enumerate}

% ---------------------------------------------------------------------
\section{第~\ref{ch:hardfault}~章 参考答案}

\begin{enumerate}
\item 典型原因包括空指针/野指针解引用、非对齐访问、执行非法函数指针、向无效地址写入等；对应代码示例可在“制造 HardFault”实验中运行。
\item 复位会丢失现场，难以定位；应先保存 LR/PC、栈指针和故障状态寄存器，再决定是离线分析还是复位恢复。
\item CFSR 说明故障类型，BFAR/MMFAR 给出出错地址，LR/PC 帮助定位执行到哪条指令；非法函数指针跳转通常表现为 PC 跳到无意义地址。
\item 非对齐访问是否触发故障与指令类型及配置有关；Cortex-M 可在控制寄存器中开启非对齐访问陷阱，便于尽早发现错误。
\item SEGGER 打印的 R0-R3、R12 是异常前函数参数与临时值，LR/PC 是调用现场，PSR 保存条件标志；结合 \texttt{.map} 文件即可反查故障函数。
\end{enumerate}

% ---------------------------------------------------------------------
\section{第~\ref{ch:linux}~章 参考答案}

\begin{enumerate}
\item \texttt{grep -rni "error" src/ --include="*.c"}
\item \texttt{ps aux | grep python} 找出 python 进程；\texttt{awk '\{print \$2\}'} 提取 PID；\texttt{xargs kill} 把标准输入逐行作为参数传给 \texttt{kill}。
\item 核心逻辑：
\begin{lstlisting}
find /tmp/logs -name "*.log" -mtime +7 -size +1M -exec rm -f {} \;
\end{lstlisting}
\item Makefile 的规则形如“目标:依赖 + 命令”；只要依赖比目标新就会执行命令。头文件也写成依赖后，修改头文件能触发重新编译。
\item \texttt{cmake -S . -B build} 用源码目录 \texttt{.} 生成构建目录 \texttt{build}；再执行 \texttt{cmake --build build} 即可编译运行 \texttt{./build/app}。
\item \texttt{-Wall -Wextra} 打开常见与额外警告，\texttt{-g} 生成调试符号；没有 \texttt{-g} 时调试器无法显示变量与源码行。
\end{enumerate}